% Options for packages loaded elsewhere
\PassOptionsToPackage{unicode}{hyperref}
\PassOptionsToPackage{hyphens}{url}
\documentclass[
  10pt,oneside]{article}
\usepackage{xcolor}
\usepackage[top=1cm,bottom=1.3cm,left=1.3cm,right=1.3cm,margin=1in]{geometry}
\usepackage{amsmath,amssymb}
\setcounter{secnumdepth}{-\maxdimen} % remove section numbering
\usepackage{iftex}
\ifPDFTeX
  \usepackage[T1]{fontenc}
  \usepackage[utf8]{inputenc}
  \usepackage{textcomp} % provide euro and other symbols
\else % if luatex or xetex
  \usepackage{unicode-math} % this also loads fontspec
  \defaultfontfeatures{Scale=MatchLowercase}
  \defaultfontfeatures[\rmfamily]{Ligatures=TeX,Scale=1}
\fi
\usepackage{lmodern}
\ifPDFTeX\else
  % xetex/luatex font selection
\fi
% Use upquote if available, for straight quotes in verbatim environments
\IfFileExists{upquote.sty}{\usepackage{upquote}}{}
\IfFileExists{microtype.sty}{% use microtype if available
  \usepackage[]{microtype}
  \UseMicrotypeSet[protrusion]{basicmath} % disable protrusion for tt fonts
}{}
\makeatletter
\@ifundefined{KOMAClassName}{% if non-KOMA class
  \IfFileExists{parskip.sty}{%
    \usepackage{parskip}
  }{% else
    \setlength{\parindent}{0pt}
    \setlength{\parskip}{6pt plus 2pt minus 1pt}}
}{% if KOMA class
  \KOMAoptions{parskip=half}}
\makeatother
\usepackage{longtable,booktabs,array}
\usepackage{calc} % for calculating minipage widths
% Correct order of tables after \paragraph or \subparagraph
\usepackage{etoolbox}
\makeatletter
\patchcmd\longtable{\par}{\if@noskipsec\mbox{}\fi\par}{}{}
\makeatother
% Allow footnotes in longtable head/foot
\IfFileExists{footnotehyper.sty}{\usepackage{footnotehyper}}{\usepackage{footnote}}
\makesavenoteenv{longtable}
\usepackage{graphicx}
\makeatletter
\newsavebox\pandoc@box
\newcommand*\pandocbounded[1]{% scales image to fit in text height/width
  \sbox\pandoc@box{#1}%
  \Gscale@div\@tempa{\textheight}{\dimexpr\ht\pandoc@box+\dp\pandoc@box\relax}%
  \Gscale@div\@tempb{\linewidth}{\wd\pandoc@box}%
  \ifdim\@tempb\p@<\@tempa\p@\let\@tempa\@tempb\fi% select the smaller of both
  \ifdim\@tempa\p@<\p@\scalebox{\@tempa}{\usebox\pandoc@box}%
  \else\usebox{\pandoc@box}%
  \fi%
}
% Set default figure placement to htbp
\def\fps@figure{htbp}
\makeatother
\setlength{\emergencystretch}{3em} % prevent overfull lines
\providecommand{\tightlist}{%
  \setlength{\itemsep}{0pt}\setlength{\parskip}{0pt}}
\usepackage{bookmark}
\IfFileExists{xurl.sty}{\usepackage{xurl}}{} % add URL line breaks if available
\urlstyle{same}
\hypersetup{
  pdfauthor={Helena Regina Salomé D'Espindula},
  hidelinks,
  pdfcreator={LaTeX via pandoc}}

\author{Helena Regina Salomé D'Espindula}
\date{2026-02-17}

\begin{document}

{
\setcounter{tocdepth}{2}
\tableofcontents
}
\section{Classificação de tecidos naturais --- panos e malhas}\label{classificauxe7uxe3o-de-tecidos-naturais-panos-e-malhas}

Tabela de referência para identificação técnica de tecidos comuns no mercado, com foco em fibras naturais (principalmente algodão e linho 100\%).

\begin{longtable}[]{@{}
  >{\raggedright\arraybackslash}p{(\linewidth - 14\tabcolsep) * \real{0.1758}}
  >{\raggedright\arraybackslash}p{(\linewidth - 14\tabcolsep) * \real{0.0659}}
  >{\raggedright\arraybackslash}p{(\linewidth - 14\tabcolsep) * \real{0.0769}}
  >{\raggedright\arraybackslash}p{(\linewidth - 14\tabcolsep) * \real{0.2088}}
  >{\raggedright\arraybackslash}p{(\linewidth - 14\tabcolsep) * \real{0.0549}}
  >{\raggedright\arraybackslash}p{(\linewidth - 14\tabcolsep) * \real{0.1209}}
  >{\raggedright\arraybackslash}p{(\linewidth - 14\tabcolsep) * \real{0.1648}}
  >{\raggedright\arraybackslash}p{(\linewidth - 14\tabcolsep) * \real{0.1319}}@{}}
\toprule\noalign{}
\begin{minipage}[b]{\linewidth}\raggedright
Nome comercial
\end{minipage} & \begin{minipage}[b]{\linewidth}\raggedright
Tipo
\end{minipage} & \begin{minipage}[b]{\linewidth}\raggedright
Fibra
\end{minipage} & \begin{minipage}[b]{\linewidth}\raggedright
Trama / Estrutura
\end{minipage} & \begin{minipage}[b]{\linewidth}\raggedright
Fio
\end{minipage} & \begin{minipage}[b]{\linewidth}\raggedright
Densidade
\end{minipage} & \begin{minipage}[b]{\linewidth}\raggedright
Transparência
\end{minipage} & \begin{minipage}[b]{\linewidth}\raggedright
Uso típico
\end{minipage} \\
\midrule\noalign{}
\endhead
\bottomrule\noalign{}
\endlastfoot
Tricoline & Pano & Algodão & Plano (tafetá) & Médio & Alta & Não & Camisas, vestidos, costura geral \\
Cambraia & Pano & Algodão & Plano & Fino & Média & Leve & Roupas leves, bebê, verão \\
Voil & Pano & Algodão & Plano (alta torção) & Muito fino & Baixa & Alta & Cortinas, blusas \\
Percal & Pano & Algodão & Plano & Fino & Muito alta & Não & Lençóis, fronhas \\
Popeline & Pano & Algodão & Plano & Fino--médio & Alta & Não & Camisas estruturadas \\
Chita & Pano & Algodão & Plano & Médio & Média & Não & Uso doméstico, decoração \\
Morim & Pano & Algodão & Plano & Médio & Média & Não & Forros, uniformes \\
Sarja leve & Pano & Algodão & Sarja & Médio & Média-alta & Não & Saias, calças \\
Brim & Pano & Algodão & Sarja & Médio--grosso & Alta & Não & Trabalho, aventais \\
Jeans (denim) & Pano & Algodão & Sarja & Grosso & Alta & Não & Calças, jaquetas \\
Piquet (tecido) & Pano & Algodão & Relevo composto & Médio & Média & Não & Colchas, polos (tecido) \\
Favo (tecido) & Pano & Algodão & Estrutural 3D & Médio & Média & Não & Toalhas, colchas \\
Panamá & Pano & Algodão / Linho & Cesta & Médio & Média & Não & Mesa, decoração \\
Canvas / Lona & Pano & Algodão & Cesta ou sarja & Grosso & Alta & Não & Bolsas, capas \\
Linho leve & Pano & Linho & Plano & Médio & Média & Leve & Camisas, vestidos \\
Linho pesado & Pano & Linho & Plano & Grosso & Alta & Não & Mesa, cortinas \\
Linho sarjado & Pano & Linho & Sarja & Médio & Média-alta & Não & Alfaiataria leve \\
Jersey (meia-malha) & Malha & Algodão & Malha simples & Fino--médio & Média & Baixa & Camisetas \\
Ribana & Malha & Algodão & Canelada & Médio & Alta & Não & Punhos, golas, cós \\
Interlock & Malha & Algodão & Malha dupla & Médio & Alta & Não & Bebê, pijamas \\
Malha canelada & Malha & Algodão & Malha simples & Médio & Média & Baixa & Vestidos ajustados \\
Piquet (malha) & Malha & Algodão & Malha estruturada & Médio & Média & Não & Camisa polo \\
Waffle / favo (malha) & Malha & Algodão & Malha 3D & Médio & Média & Não & Roupões, mantas \\
Malha de linho & Malha & Linho & Jersey & Médio & Média & Leve & Peças de verão \\
\end{longtable}

\textbf{Nota técnica:}\\
O nome comercial do tecido é consequência da combinação entre fibra, tipo estrutural (pano ou malha), densidade, fio e uso histórico. Para decisões de costura consciente, priorizar sempre a identificação do tipo (pano × malha) e da densidade antes do nome.

\newpage

\section{Blusa}\label{blusa}

\section{Projeto: Blusa Esportiva para Pilates}\label{projeto-blusa-esportiva-para-pilates}

Data ideia: 21.01.2026

\textbf{Versão 1.0}
Data início modelagem:
Data início tecido:
Data finalizada:

\begin{itemize}
\tightlist
\item
  Modelagem simples com pala frontal e gusset
\item
  Tecido plano leve -- foco em conforto, mobilidade e aprendizado
\end{itemize}

\begin{center}\rule{0.5\linewidth}{0.5pt}\end{center}

\subsection{1. Objetivo do projeto}\label{objetivo-do-projeto}

Desenvolver uma blusa esportiva leve, confortável e funcional para prática de Pilates, utilizando tecido plano natural (tricoline leve 100\% algodão), com foco em:

\begin{itemize}
\tightlist
\item
  conforto físico e psicológico
\item
  liberdade de movimento
\item
  ausência de compressão corporal
\item
  acabamento interno confortável
\item
  aprendizado de modelagem e construção
\item
  soluções inspiradas em alfaiataria histórica e funcional
\end{itemize}

O projeto \textbf{não visa performance atlética extrema}, e sim uso real, repetido e confortável.

\begin{center}\rule{0.5\linewidth}{0.5pt}\end{center}

\subsection{2. Premissas do projeto}\label{premissas-do-projeto}

\begin{itemize}
\tightlist
\item
  Tecido \textbf{não elástico}
\item
  Peça \textbf{não justa}
\item
  Mobilidade resolvida por \textbf{construção}, não por compressão
\item
  Prioridade para:

  \begin{itemize}
  \tightlist
  \item
    estabilidade
  \item
    previsibilidade
  \item
    facilidade de ajuste
  \end{itemize}
\item
  Técnicas compatíveis com:

  \begin{itemize}
  \tightlist
  \item
    iniciante em modelagem
  \item
    costura doméstica
  \item
    ausência de overloque
  \end{itemize}
\end{itemize}

\begin{center}\rule{0.5\linewidth}{0.5pt}\end{center}

\subsection{3. Especificações gerais da peça}\label{especificauxe7uxf5es-gerais-da-peuxe7a}

\subsubsection{Tipo de peça}\label{tipo-de-peuxe7a}

Blusa esportiva leve, manga curta, comprimento médio.

\subsubsection{Uso previsto}\label{uso-previsto}

Pilates (atividade física leve a moderada).

\subsubsection{Silhueta}\label{silhueta}

\begin{itemize}
\tightlist
\item
  solta
\item
  organizada
\item
  sem efeito ``saco de batata''
\item
  sem marcação de abdômen ou busto
\end{itemize}

\subsubsection{Estrutura}\label{estrutura}

\begin{itemize}
\tightlist
\item
  frente com \textbf{pala horizontal}
\item
  costas em \textbf{duas peças com costura central}
\item
  mangas curtas simples
\item
  \textbf{gusset discreto na axila}
\item
  sem pences
\item
  sem elástico
\item
  sem recortes complexos
\end{itemize}

\begin{center}\rule{0.5\linewidth}{0.5pt}\end{center}

\subsection{4. Tecido}\label{tecido}

\subsubsection{Tecido principal}\label{tecido-principal}

\begin{itemize}
\tightlist
\item
  Tricoline leve 100\% algodão
\end{itemize}

\subsubsection{Justificativa}\label{justificativa}

\begin{itemize}
\tightlist
\item
  respirável
\item
  absorvente
\item
  confortável ao toque
\item
  didaticamente adequado
\item
  tecnicamente aceitável para Pilates
\end{itemize}

\subsubsection{Limitações conhecidas}\label{limitauxe7uxf5es-conhecidas}

\begin{itemize}
\tightlist
\item
  seca mais lentamente que tecidos sintéticos
\item
  pode ficar levemente transparente com suor
\end{itemize}

Essas limitações são consideradas aceitáveis no contexto do projeto.

\begin{center}\rule{0.5\linewidth}{0.5pt}\end{center}

\subsection{5. Medidas necessárias (todas em cm)}\label{medidas-necessuxe1rias-todas-em-cm}

Somente medidas essenciais:

\begin{enumerate}
\def\labelenumi{\arabic{enumi}.}
\tightlist
\item
  Contorno de busto
\item
  Largura de ombros
\item
  Altura do ombro até comprimento desejado da blusa
\item
  Circunferência do braço alto
\item
  Comprimento da manga curta
\end{enumerate}

\begin{center}\rule{0.5\linewidth}{0.5pt}\end{center}

\subsection{6. Folgas recomendadas}\label{folgas-recomendadas}

\subsubsection{Folga de conforto (iniciante)}\label{folga-de-conforto-iniciante}

\begin{itemize}
\tightlist
\item
  Busto: \textbf{+8 a +12 cm} no total
\item
  Braço: \textbf{+4 a +6 cm}
\item
  Comprimento: não justo ao quadril
\end{itemize}

A folga é distribuída, não concentrada.

\begin{center}\rule{0.5\linewidth}{0.5pt}\end{center}

\subsection{7. Estrutura do molde}\label{estrutura-do-molde}

\subsubsection{7.1 Frente}\label{frente}

\paragraph{Base}\label{base}

\begin{itemize}
\tightlist
\item
  Retângulo
\item
  Altura = ombro até comprimento final
\item
  Largura = (busto ÷ 2) + folga
\end{itemize}

\paragraph{Decote}\label{decote}

\begin{itemize}
\tightlist
\item
  Raso
\item
  Confortável
\item
  Sem aprofundamento excessivo
\end{itemize}

\paragraph{Pala frontal}\label{pala-frontal}

\begin{itemize}
\tightlist
\item
  Linha horizontal marcada:

  \begin{itemize}
  \tightlist
  \item
    logo \textbf{acima do ponto mais alto do busto}
  \end{itemize}
\item
  Divide a frente em:

  \begin{itemize}
  \tightlist
  \item
    pala superior
  \item
    frente inferior
  \end{itemize}
\end{itemize}

A pala:
- organiza o tecido
- estabiliza a parte superior
- evita efeito ``saco''

A pala será copiada como peça separada.

\begin{center}\rule{0.5\linewidth}{0.5pt}\end{center}

\subsubsection{7.2 Costas}\label{costas}

\begin{itemize}
\tightlist
\item
  Mesmo retângulo base da frente
\item
  Divididas em \textbf{duas peças}
\item
  Costura central nas costas (não cortar na dobra)
\end{itemize}

Funções da costura central:
- criar eixo visual
- ajudar no assentamento do tecido
- permitir ajustes futuros
- evitar excesso de pano solto

Decote das costas:
- mais alto que o da frente
- simples

\begin{center}\rule{0.5\linewidth}{0.5pt}\end{center}

\subsubsection{7.3 Manga}\label{manga}

\begin{itemize}
\tightlist
\item
  Manga curta simples
\item
  Retângulo no fio reto
\end{itemize}

Dimensões:
- Comprimento: \textbf{5 a 8 cm}
- Largura: circunferência do braço + folga

A mobilidade do braço será resolvida pelo gusset, não pela cava.

\begin{center}\rule{0.5\linewidth}{0.5pt}\end{center}

\subsubsection{7.4 Gusset (axila)}\label{gusset-axila}

\begin{itemize}
\tightlist
\item
  Função:

  \begin{itemize}
  \tightlist
  \item
    liberar movimento
  \item
    reduzir tensão
  \item
    melhorar conforto térmico
  \item
    evitar subida da peça
  \end{itemize}
\end{itemize}

Formato recomendado:
- losango curto \textbf{ou}
- formato ``folha'' simples (mais discreto)

Dimensões iniciais:
- largura: \textbf{6 a 8 cm}
- comprimento: \textbf{8 a 10 cm}

Observação:
Começar menor é mais seguro; é mais fácil aumentar do que reduzir.

\begin{center}\rule{0.5\linewidth}{0.5pt}\end{center}

\subsection{8. Margens de costura (aprendizado)}\label{margens-de-costura-aprendizado}

\begin{itemize}
\tightlist
\item
  Laterais: \textbf{2 a 3 cm}
\item
  Ombros: \textbf{1,5 a 2 cm}
\item
  Costura central das costas: \textbf{2 cm}
\item
  Pala: \textbf{1,5 cm}
\item
  Manga: \textbf{1,5 cm}
\item
  Gusset: \textbf{1 cm}
\item
  Barra: \textbf{3 a 4 cm}
\end{itemize}

Margens maiores permitem ajustes sem comprometer a peça.

\begin{center}\rule{0.5\linewidth}{0.5pt}\end{center}

\subsection{9. Acabamentos internos}\label{acabamentos-internos}

\subsubsection{Prioridade}\label{prioridade}

\begin{itemize}
\tightlist
\item
  conforto
\item
  ausência de aspereza
\item
  baixo volume
\end{itemize}

\subsubsection{Técnicas recomendadas}\label{tuxe9cnicas-recomendadas}

\begin{itemize}
\tightlist
\item
  \textbf{costura francesa}:

  \begin{itemize}
  \tightlist
  \item
    laterais
  \item
    ombros
  \item
    mangas
  \end{itemize}
\item
  \textbf{rebatimento interno discreto} da costura francesa:

  \begin{itemize}
  \tightlist
  \item
    quando houver projeção excessiva
  \item
    feito com ponto longo
  \item
    costurando apenas na margem interna
  \end{itemize}
\end{itemize}

\subsubsection{Evitar}\label{evitar}

\begin{itemize}
\tightlist
\item
  zigue-zague exposto
\item
  viés grosso na axila
\item
  costuras duras ou espessas
\end{itemize}

\begin{center}\rule{0.5\linewidth}{0.5pt}\end{center}

\subsection{10. Acabamentos de borda}\label{acabamentos-de-borda}

\subsubsection{Pala frontal (borda inferior)}\label{pala-frontal-borda-inferior}

\begin{itemize}
\tightlist
\item
  bainha estreita (0,5 + 0,5 cm)
\item
  ou debrum leve
\end{itemize}

\subsubsection{Barras}\label{barras}

\begin{itemize}
\tightlist
\item
  barra simples ou dupla
\item
  largura suficiente para peso leve
\item
  sem rigidez
\end{itemize}

\begin{center}\rule{0.5\linewidth}{0.5pt}\end{center}

\subsection{11. Ordem geral de montagem (resumo)}\label{ordem-geral-de-montagem-resumo}

\begin{enumerate}
\def\labelenumi{\arabic{enumi}.}
\tightlist
\item
  Preparar a pala (acabar borda inferior)
\item
  Aplicar pala à frente
\item
  Costurar centro das costas
\item
  Costurar ombros
\item
  Preparar mangas
\item
  Inserir mangas
\item
  Inserir gusset
\item
  Fechar laterais
\item
  Provar e ajustar
\item
  Fazer barras
\end{enumerate}

\begin{center}\rule{0.5\linewidth}{0.5pt}\end{center}

\subsection{12. Critérios de sucesso do projeto}\label{crituxe9rios-de-sucesso-do-projeto}

A peça final deve:

\begin{itemize}
\tightlist
\item
  permitir levantar os braços sem subir excessivamente
\item
  não apertar busto, abdômen ou axila
\item
  não marcar gordura corporal
\item
  não causar desconforto interno
\item
  permitir deitar, torcer e mover o tronco
\item
  manter leitura visual organizada
\item
  ser usável repetidamente
\end{itemize}

\begin{center}\rule{0.5\linewidth}{0.5pt}\end{center}

\subsection{13. Observação sobre uso consciente}\label{observauxe7uxe3o-sobre-uso-consciente}

Quando possível:
- aproveitar sobras para gussets ou viés
- priorizar ajustes antes de cortes definitivos
- pensar em manutenção futura

Nesta fase, o \textbf{aprendizado técnico tem prioridade} sobre otimização extrema de material.

\begin{center}\rule{0.5\linewidth}{0.5pt}\end{center}

\subsection{14. Notas finais}\label{notas-finais}

Este molde não visa perfeição estética imediata, mas:

\begin{itemize}
\tightlist
\item
  funcionalidade real
\item
  compreensão de construção
\item
  base reutilizável para versões futuras
\item
  evolução gradual de complexidade
\end{itemize}

Projeto pensado para durar, ser ajustado e reaproveitado.

\section{Saias}\label{saias}

\section{1. Identificação do Projeto}\label{identificauxe7uxe3o-do-projeto}

\begin{itemize}
\tightlist
\item
  \textbf{Tipo de projeto:} Reforma e adaptação de peça esportiva existente\\
\item
  \textbf{Finalidade:} Uso em Pilates em aparelho\\
\item
  \textbf{Usuária:} Corpo real (plus size, barriga e quadril pronunciados)\\
\item
  \textbf{Etapa de aprendizagem:} Modelagem consciente e funcional\\
\item
  \textbf{Manequim:} Não utilizado (prova no corpo e em superfície plana)
\end{itemize}

\begin{center}\rule{0.5\linewidth}{0.5pt}\end{center}

\section{2. Peça Original -- Especificações}\label{peuxe7a-original-especificauxe7uxf5es}

\subsection{2.1 Descrição da peça original}\label{descriuxe7uxe3o-da-peuxe7a-original}

\begin{itemize}
\tightlist
\item
  \textbf{Tipo:} Short-saia esportivo
\item
  \textbf{Uso:} Atividade física
\item
  \textbf{Tecido:} Malha esportiva (provável poliamida + elastano)
\item
  \textbf{Cós:} Faixa elástica contínua
\item
  \textbf{Cor:} Roxo quase neon
\item
  \textbf{Observação funcional:}\\
  Saia curta, sobe facilmente em movimento, sensação de exposição excessiva (body shaming).
\end{itemize}

\subsection{2.2 Decisão de projeto}\label{decisuxe3o-de-projeto}

\begin{itemize}
\tightlist
\item
  \textbf{Não remover a saia original nesta etapa}
\item
  Motivos:

  \begin{itemize}
  \tightlist
  \item
    Manter reversibilidade
  \item
    Evitar intervenção irreversível precoce
  \item
    Avaliar comportamento real após uso
  \end{itemize}
\end{itemize}

\begin{center}\rule{0.5\linewidth}{0.5pt}\end{center}

\section{3. Materiais Disponíveis}\label{materiais-disponuxedveis}

\subsection{3.1 Tecidos}\label{tecidos}

\begin{itemize}
\item
  \textbf{Tricoline 100\% algodão -- preto (peso médio):}\\
  Dimensão aproximada: 3,0 m × 1,40 m\\
  Uso: estrutura principal da nova saia
\item
  \textbf{Tricoline 100\% algodão -- lilás (bem leve):}\\
  Dimensão aproximada: 1,0 m × 1,40 m\\
  Uso futuro: forro parcial inferior / efeito visual (não estrutural)
\end{itemize}

\subsection{3.2 Linhas e insumos}\label{linhas-e-insumos}

\begin{itemize}
\tightlist
\item
  Linha de poliéster preta
\item
  Ferro de passar (fundamental para bainha e controle de caimento)
\end{itemize}

\begin{center}\rule{0.5\linewidth}{0.5pt}\end{center}

\section{4. Medidas Corporais e da Peça Original}\label{medidas-corporais-e-da-peuxe7a-original}

\subsection{4.1 Medidas do short-saia original}\label{medidas-do-short-saia-original}

\begin{itemize}
\tightlist
\item
  Circunferência do cós (fora do corpo): \textbf{79 cm}
\item
  Altura do cós: \textbf{4,5 cm}
\item
  Base do cós até fim do short:

  \begin{itemize}
  \tightlist
  \item
    Fora do corpo: \textbf{28 cm}
  \item
    No corpo: \textbf{37 cm}
  \end{itemize}
\item
  Topo do cós até o joelho (no corpo): \textbf{54,5 cm}
\end{itemize}

\subsection{4.2 Circunferência funcional adotada}\label{circunferuxeancia-funcional-adotada}

\begin{itemize}
\tightlist
\item
  Ajuste de +12\% sobre a medida fora do corpo
\item
  Circunferência funcional de trabalho (A): \textbf{88,5 cm}
\end{itemize}

\begin{center}\rule{0.5\linewidth}{0.5pt}\end{center}

\section{5. Proposta de Intervenção}\label{proposta-de-intervenuxe7uxe3o}

\subsection{5.1 Conceito}\label{conceito}

\begin{itemize}
\tightlist
\item
  \textbf{Saia envelope tipo ballet (romantique funcional)}
\item
  Comprimento abaixo do joelho
\item
  Abertura lateral anterior
\item
  Presa somente na \textbf{base do cós}
\item
  Pensada para:

  \begin{itemize}
  \tightlist
  \item
    amplitude de movimento
  \item
    estabilidade
  \item
    cobertura corporal
  \item
    conforto psicológico durante a prática
  \end{itemize}
\end{itemize}

\subsection{5.2 Justificativa do modelo}\label{justificativa-do-modelo}

\begin{itemize}
\tightlist
\item
  Envelope permite abertura de pernas sem travar
\item
  Comprimento longo reduz exposição em movimento
\item
  Fixação no cós evita ajustes durante a aula
\item
  Adequado para Pilates em aparelho
\end{itemize}

\begin{center}\rule{0.5\linewidth}{0.5pt}\end{center}

\section{6. Modelagem -- Saia Preta (Estrutura)}\label{modelagem-saia-preta-estrutura}

\subsection{6.1 Tipo de molde}\label{tipo-de-molde}

\begin{itemize}
\tightlist
\item
  Peça única
\item
  Assimétrica
\item
  Corte único (não espelhado)
\item
  Fio reto central
\end{itemize}

\subsection{6.2 Medidas do molde}\label{medidas-do-molde}

\subsubsection{Comprimento (linha central / fio)}\label{comprimento-linha-central-fio}

\begin{itemize}
\tightlist
\item
  Comprimento útil (base do cós até abaixo do joelho): \textbf{55 cm}
\item
  Bainha com peso: \textbf{5 cm}
\item
  Margem superior: \textbf{1,5 cm}
\item
  \textbf{Comprimento total:} \textbf{62 cm}
\end{itemize}

\subsubsection{Largura do topo}\label{largura-do-topo}

\begin{itemize}
\tightlist
\item
  ½ circunferência funcional: \textbf{44,25 cm}
\item
  Sobreposição frontal: \textbf{12 cm}
\item
  \textbf{Topo total:} \textbf{56,5 cm}
\end{itemize}

\subsubsection{Largura máxima da base (antes da curva)}\label{largura-muxe1xima-da-base-antes-da-curva}

\begin{itemize}
\tightlist
\item
  Abertura adicional: \textbf{+30 cm}
\item
  \textbf{Base teórica:} \textbf{87 cm}
\end{itemize}

\begin{center}\rule{0.5\linewidth}{0.5pt}\end{center}

\subsection{6.3 Geometria da barra}\label{geometria-da-barra}

\begin{itemize}
\tightlist
\item
  Centro da peça: +3 a +4 cm mais longo
\item
  Laterais: 3 a 4 cm mais curtas
\item
  Curva contínua, sem ângulos
\item
  Bainha marcada a 5 cm
\end{itemize}

\begin{center}\rule{0.5\linewidth}{0.5pt}\end{center}

\subsection{6.4 Marcas obrigatórias no molde}\label{marcas-obrigatuxf3rias-no-molde}

\begin{itemize}
\tightlist
\item
  Linha do fio com setas duplas
\item
  Centro da peça
\item
  Pontos de início e fim da fixação no cós (±14 cm do centro)
\item
  Linha de referência do fim do short interno (37 cm do topo)
\item
  Identificação completa do modelo
\item
  Indicação clara: \textbf{margem de costura não inclusa}
\item
  Observação: \textbf{não cortar no viés}
\end{itemize}

\begin{center}\rule{0.5\linewidth}{0.5pt}\end{center}

\section{7. Esboço de Melhoria Futura -- Forro Lilás}\label{esbouxe7o-de-melhoria-futura-forro-liluxe1s}

\subsection{7.1 Conceito}\label{conceito-1}

\begin{itemize}
\tightlist
\item
  Forro parcial inferior
\item
  Função:

  \begin{itemize}
  \tightlist
  \item
    estabilização fina da barra
  \item
    efeito visual
  \end{itemize}
\item
  Não estrutural
\item
  Não sobe até o cós
\end{itemize}

\subsection{7.2 Esboço no molde}\label{esbouxe7o-no-molde}

\begin{itemize}
\tightlist
\item
  Linha tracejada paralela à barra
\item
  Altura sugerida: \textbf{12 a 18 cm} (referência adotada: 15 cm)
\item
  Anotações explícitas:

  \begin{itemize}
  \tightlist
  \item
    ``Forro parcial inferior''
  \item
    ``Aplicação posterior''
  \item
    ``Não estrutural''
  \end{itemize}
\end{itemize}

\begin{center}\rule{0.5\linewidth}{0.5pt}\end{center}

\section{8. Montagem e Costura -- Sequência Planejada}\label{montagem-e-costura-sequuxeancia-planejada}

\subsection{8.1 Ordem de execução}\label{ordem-de-execuuxe7uxe3o}

\begin{enumerate}
\def\labelenumi{\arabic{enumi}.}
\tightlist
\item
  Desenho e corte do molde em tricoline preto
\item
  Costura da saia preta à \textbf{base do cós original}
\item
  Execução da bainha de 5 cm
\item
  Costura dupla paralela na bainha (peso)
\item
  Prova no corpo em uso real (Pilates em aparelho)
\item
  Avaliação de necessidade de forro lilás futuro
\end{enumerate}

\subsection{8.2 Considerações de costura}\label{considerauxe7uxf5es-de-costura}

\begin{itemize}
\tightlist
\item
  Não costurar laterais
\item
  Não atravessar o elástico do cós
\item
  Manter costuras estruturais e estéticas separadas
\item
  Priorizar conforto em posição sentada e em extensão
\end{itemize}

\begin{center}\rule{0.5\linewidth}{0.5pt}\end{center}

\section{9. Critérios de Avaliação Pós-Uso}\label{crituxe9rios-de-avaliauxe7uxe3o-puxf3s-uso}

\begin{itemize}
\tightlist
\item
  A saia sobe em movimento?
\item
  Há giro excessivo?
\item
  A barra retorna rapidamente após abertura?
\item
  Há desconforto térmico?
\item
  A saia original interfere no caimento?
\end{itemize}

Somente após essas observações será avaliada:
- remoção da saia original
- instalação do forro lilás

\begin{center}\rule{0.5\linewidth}{0.5pt}\end{center}

\section{10. Observações Finais}\label{observauxe7uxf5es-finais}

Este projeto foi concebido como exercício de:
- modelagem consciente
- adaptação ao corpo real
- priorização de uso e conforto
- documentação para aprendizagem contínua

Decisões foram tomadas visando reversibilidade e evolução progressiva do projeto.

\begin{center}\rule{0.5\linewidth}{0.5pt}\end{center}

\section{Vestidos}\label{vestidos}

\vspace*{-1cm}

\begin{center}
\normalsize\bfseries PROJETO:
\end{center}

\vspace{0.8cm}

\begin{center}
\setlength{\tabcolsep}{18pt}
\renewcommand{\arraystretch}{1.15}
\begin{tabular}{cc}

\begin{minipage}[t]{0.47\textwidth}
\normalsize

\textbf{Identificação}

\vspace{0.3cm}

Projeto: \rule{6cm}{0.4pt}  

Data de início: \rule{4cm}{0.4pt}  

Molde base: \rule{6cm}{0.4pt}  

Referência (época/estilo): \rule{3cm}{0.4pt}

\rule{6cm}{0.4pt}

\vspace{1cm}

\end{minipage}
&
\begin{minipage}[t]{0.47\textwidth}
\normalsize

\textbf{Materiais}

\vspace{0.3cm}

Tecido principal: \rule{5cm}{0.4pt}  

Composição: \rule{6cm}{0.4pt}  

Metragem estimada: \rule{4cm}{0.4pt}  

Aviamentos: \rule{6cm}{0.4pt}

\end{minipage}

\\[2cm]

\multicolumn{2}{c}{
\begin{minipage}[t]{0.95\textwidth}
\normalsize

\textbf{Objetivo da peça}

\vspace{0.3cm}

Escreva aqui o propósito da peça, contexto de uso, referências estéticas e intenção de silhueta.

\vspace{0.8cm}

\textbf{Planejamento de construção}

\vspace{0.3cm}

Descreva a sequência planejada de montagem, pontos críticos da modelagem e acabamentos pretendidos.

\vspace{0.8cm}

\textbf{Ajustes realizados durante a execução}

\vspace{0.3cm}

\rule{15cm}{0.4pt}  

\rule{15cm}{0.4pt}  

\rule{15cm}{0.4pt}  

\rule{15cm}{0.4pt}

\end{minipage}
}

\end{tabular}
\end{center}

\end{document}
